% !TeX root = ../../Book1.tex
% 纲要标题：### 17.2 自由固有作用的商定理
% 纲要位置：几何学计划纲要.md，第 779 行。
% 正文状态：初稿；尚未正式审核。

% 纲要内容开始
% 
% * 固有性的局部和序列判据
% * 局部横截面与局部乘积
% * 商的 Hausdorff 性和光滑图册
% * 商映射的泛性质
% * 缺少自由或固有条件的反例
% 
% 纲要内容结束

\section{自由固有作用的商定理}
\label{g1:sec:1702}

把同一轨道上的点合并以后，横向方向应当留下来，沿轨道运动的方向则应被忘掉。自由性保证各条轨道具有同一个维数，但它不能阻止一条轨道反复回到很小的区域。为了让横向坐标真正成为商空间的坐标，还需要控制这种返回。本节的固有性正是一个同时限制出发点、到达点和群参数的条件。

\subsection{固有性的局部和序列判据}
\label{g1:subsec:1702:01}

\begin{definition}[固有作用]\label{g1:def:proper-lie-group-action}
	设有限维实 Lie 群 \(G\) 光滑左作用于有限维无边界光滑流形 \(M\)，两者均为 Hausdorff、第二可数。若映射
	\[
		\Theta:G\times M\longrightarrow M\times M,\qquad
		\Theta(g,p)=(gp,p)
	\]
	为固有映射，即每个紧集的原像都紧，则称这个作用为\term[guyouzuoyong]{固有作用}[proper action]。
\end{definition}

这里采用的是\cref{g1:def:proper-map-manifolds}的固有映射概念。需要约束的映射是 \(\Theta\)，而非单独的作用映射 \((g,p)\mapsto gp\)：后者忘掉出发点，即使直线上的平移也不满足同样的紧集原像条件。对两个子集 \(A,B\subseteq M\)，记
\[
	T(A,B)=\{g\in G:gA\cap B\ne\varnothing\}.
\]
它收集所有能够把 \(A\) 中某点送入 \(B\) 的群元素，下面称为从 \(A\) 到 \(B\) 的搬运集。这个记号中的先后次序与 \(\Theta\) 的两个输出次序不同：\(\Theta^{-1}(B\times A)\) 投影到群因子，才得到 \(T(A,B)\)。

\begin{proposition}[固有性的判据]\label{g1:prop:proper-action-criteria}
	设有限维实 Lie 群 \(G\) 光滑作用于有限维无边界光滑流形 \(M\)，且两者均为 Hausdorff、第二可数。对 \(A,B\subseteq M\)，令 \(T(A,B)=\{g\in G:gA\cap B\ne\varnothing\}\)。下列条件等价。
	\begin{equivalences}
		\item 作用固有。
		\item 对任意紧集 \(K,L\subseteq M\)，\(T(K,L)\) 在 \(G\) 中紧。
		\item 对任意 \(p,q\in M\)，存在开邻域 \(U\ni p\)、\(V\ni q\)，使 \(T(U,V)\) 在 \(G\) 中的闭包紧。
		\item 若 \(p_j\to p\)、\(g_jp_j\to q\) 在 \(M\) 中收敛，则 \((g_j)\) 有收敛子列。
	\end{equivalences}
	第二个条件也可只对 \(K=L\) 检查。
\end{proposition}
\begin{proof}
	先说明本证明使用序列的依据。流形有可数局部基，紧子集中的任意序列都有收敛子列：若无无限重复值，紧性给出一个聚点，再沿其递减局部基逐次选出子列。反过来，在第二可数空间的子集中，若每个序列都有收敛子列，则该子集紧。否则一个没有有限子覆盖的开覆盖可缩成可数子覆盖 \(O_1,O_2,\ldots\)，逐次选取 \(x_j\notin O_1\cup\cdots\cup O_j\)。任何收敛子列的极限都落在某个 \(O_k\)，从而子列最终也落在 \(O_k\)，矛盾。这使以下序列论证适用于本命题，而无须把它误用到任意拓扑空间。

	I 蕴含 II，因为 \(\Theta^{-1}(L\times K)\) 紧，而它在 \(G\) 中的投影就是 \(T(K,L)\)。反过来，设 II 成立，取紧集 \(C\subseteq M\times M\)，其两个投影分别记为 \(L,K\)。那么
	\[
		\Theta^{-1}(C)\subseteq T(K,L)\times K.
	\]
	右侧紧；\(C\) 在 Hausdorff 空间中闭，故左侧是右侧的闭子集，也紧，证明 I。若仅已知每个 \(T(K,K)\) 紧，则 \(T(K,L)\subseteq T(K\cup L,K\cup L)\)。此外 \(T(K,L)\) 总是闭的：若 \(g_j\to g\) 且 \(g_jx_j\in L\)、\(x_j\in K\)，紧性允许取 \(x_j\to x\in K\) 的子列，连续性给出 \(gx\in L\)，从而 \(g\in T(K,L)\)。流形中闭性可由序列检验，所以紧集的这个闭子集仍紧。

	II 蕴含 III，只须在 \(p,q\) 的坐标图中各取一个闭包紧的小开球。若反过来已知 III，对紧集 \(K,L\)，用相应乘积邻域 \(U\times V\) 覆盖 \(K\times L\)，再取有限子覆盖 \(U_i\times V_i\)。每个 \(g\in T(K,L)\) 都有见证点 \(x\in K\)、\(gx\in L\)；这对点落在某个 \(U_i\times V_i\)，所以
	\[
		T(K,L)\subseteq\bigcup_i\overline{T(U_i,V_i)}.
	\]
	右侧是有限个紧集之并。刚才已证左侧闭，故它紧。这也解释了为什么局部条件要允许任意一对 \(p,q\)，而不只检查同一点附近。

	I 蕴含 IV：序列 \((g_jp_j,p_j)\) 连同其极限组成紧集，固有性使全部 \((g_j,p_j)\) 落在一个紧集内，再取收敛子列即可。反之，若 IV 成立，对任意紧集 \(C\subseteq M\times M\) 及任意序列 \((g_j,p_j)\in\Theta^{-1}(C)\)，先由 \(C\) 紧取子列，使 \(p_j\to p\)、\(g_jp_j\to q\)。再由 IV 取 \(g_j\to g\) 的子列，便有 \((g_j,p_j)\to(g,p)\)。由于 \(C\) 闭且 \(\Theta\) 连续，这个极限仍在 \(\Theta^{-1}(C)\) 中。因此原像序列紧，按开头的说明即为紧集。
\end{proof}

若 \(G\) 离散，紧子集只有有限个点，于是第二个判据成为“每个紧集只被有限个群元素搬回与自身相交的位置”。这恢复了\secref{g1:sec:1003}讨论离散等距群时采用的固有性，但此时仍允许有限而非平凡的稳定子。

\begin{proposition}[固有作用的稳定子与轨道]\label{g1:prop:proper-action-orbits}
	设有限维实 Lie 群 \(G\) 固有地光滑作用于有限维无边界光滑流形 \(M\)，两者均为 Hausdorff、第二可数。则每个稳定子 \(G_p\) 紧，每条轨道 \(Gp\) 为 \(M\) 中的闭嵌入子流形，而且轨道等价关系
	\[
		R=\{(q,p)\in M\times M:q\in Gp\}
	\]
	在 \(M\times M\) 中闭。
\end{proposition}
\begin{proof}
	稳定子恰为 \(T(\{p\},\{p\})\)，因此紧。若 \((q_j,p_j)\in R\) 收敛到 \((q,p)\)，选取 \(g_j\) 使 \(q_j=g_jp_j\)。由序列判据，取子列后 \(g_j\to g\)，于是 \(q=gp\)，所以 \((q,p)\in R\)。这证明 \(R\) 闭。固定第二个变量为 \(p\)，\(Gp\) 就是连续映射 \(q\mapsto(q,p)\) 下 \(R\) 的原像，因而闭。

	为说明嵌入性，对任意紧集 \(C\subseteq M\)，轨道映射的原像 \(\alpha_p^{-1}(C)=T(\{p\},C)\) 紧。令 \(q:G\to G/G_p\)，则 \(\overline\alpha_p^{-1}(C)=q(\alpha_p^{-1}(C))\) 也紧。因此\cref{g1:prop:orbit-immersion}的单射浸入 \(\overline\alpha_p:G/G_p\to M\) 固有，由\cref{g1:prop:proper-injective-immersion}，它是闭像嵌入。这也说明在固有作用下，轨道的自然流形拓扑就是环境子空间拓扑。
\end{proof}

\begin{proposition}[紧群作用]\label{g1:prop:compact-group-action-proper}
	设紧的有限维实 Lie 群 \(G\) 光滑作用于有限维、Hausdorff、第二可数的无边界光滑流形 \(M\)。则作用固有。
\end{proposition}
\begin{proof}
	对紧集 \(C\subseteq M\times M\)，记第二个投影像为 \(K\)。\(\Theta^{-1}(C)\) 是紧集 \(G\times K\) 的闭子集，所以紧。
\end{proof}

\subsection{局部横截面与局部乘积}
\label{g1:subsec:1702:02}

在自由作用下，\cref{g1:prop:orbit-map-rank}给出单射 \(\mathfrak g\to \mathrm{T}_pM\)、\(X\mapsto X_M(p)\)，其像是轨道方向。选一个线性补空间，再用流形坐标把它实现成经过 \(p\) 的小子流形 \(S\)。这份\term[jubuhengjiemian]{局部横截面}[local transversal]的任务是为附近每条轨道提供一个代表。逆函数定理首先只保证小群参数内的代表唯一；固有性还将排除其他群元素把横截面送回自身。

\begin{lemma}[自由固有作用的局部乘积]\label{g1:lem:free-proper-local-product}
	设有限维实 Lie 群 \(G\) 自由且固有地光滑作用于 \(m\) 维无边界光滑流形 \(M\)，两者均为 Hausdorff、第二可数。令 \(r=\dim G\)。对每个 \(p\in M\)，存在经过 \(p\) 的 \(m-r\) 维嵌入子流形 \(S\)，使
	\begin{equation}\label{g1:eq:free-action-local-product}
		\Phi:G\times S\longrightarrow GS,\qquad \Phi(g,s)=gs
	\end{equation}
	为微分同胚，且饱和集 \(GS=\{gs:g\in G,\ s\in S\}\) 在 \(M\) 中开。每条与 \(S\) 相交的轨道只与 \(S\) 相交一次。
\end{lemma}
\begin{proof}
	先取 \(\mathrm{T}_pM\) 的直和分解
	\[
		\mathrm{T}_pM=\{X_M(p):X\in\mathfrak g\}\oplus E.
	\]
	在以 \(p\) 为原点的坐标图中，对坐标再作可逆线性变换，就可令最后 \(m-r\) 个坐标方向对应 \(E\)。前 \(r\) 个坐标置零，得到经过 \(p\) 的嵌入子流形 \(S_0\)，且 \(\mathrm{T}_pS_0=E\)。作用映射限制为 \(\Phi:G\times S_0\to M\)，其在 \((e,p)\) 的微分是
	\[
		\mathrm{d}\Phi_{(e,p)}(X,v)=X_M(p)+v.
	\]
	自由性和上述直和使它为同构。由\cref{g1:thm:inverse-function}，可以选取单位元开邻域 \(U\subseteq G\) 与 \(p\) 在 \(S_0\) 中的开邻域 \(S_1\)，使 \(\Phi|_{U\times S_1}\) 为到 \(M\) 中一个开集的微分同胚。此时 \((e,s)\)、\(s\in S_1\)，也都具有可逆微分。

	下一步是控制 \(U\) 外的群元素。我们断言，存在 \(p\) 在 \(M\) 中的开邻域 \(V\)，使
	\begin{equation}\label{g1:eq:proper-action-no-return}
		gV\cap V\ne\varnothing\quad\Longrightarrow\quad g\in U.
	\end{equation}
	若否，在 \(p\) 的一列递减坐标邻域 \(V_j\) 中，可选 \(s_j,t_j\to p\) 及 \(g_j\notin U\)，使 \(g_js_j=t_j\)。固有性的序列判据给出收敛子列 \(g_j\to g\)。取极限得 \(gp=p\)，自由性于是迫使 \(g=e\)。由于 \(U\) 是 \(e\) 的开邻域，这与 \(g_j\notin U\) 矛盾，证明断言。

	取 \(S=S_1\cap V\)，必要时再缩小为其中经过 \(p\) 的坐标邻域。如果 \(gs=t\)，其中 \(s,t\in S\)，则\eqref{g1:eq:proper-action-no-return}说明 \(g\in U\)。现在
	\[
		\Phi(g,s)=t=\Phi(e,t),
	\]
	而 \(\Phi\) 在 \(U\times S_1\) 上单射，故 \(g=e\)、\(s=t\)。若一般地 \(gs=ht\)，则 \(h^{-1}gs=t\)，由同一结论得 \(g=h\)、\(s=t\)。因此整个 \(G\times S\) 上的 \(\Phi\) 单射，而不只是小群邻域上的单射。

	还须说明这个全局于群因子的映射具有光滑逆。固定 \(g\in G\)，由
	\[
		\Phi(gk,s)=g\Phi(k,s)
	\]
	可知 \(\mathrm{d}\Phi_{(g,s)}\) 是 \(\mathrm{d}\Phi_{(e,s)}\) 与两侧平移微分同构的复合，所以处处可逆。逆函数定理使 \(\Phi\) 为局部微分同胚，因而为开映射，其像 \(GS\) 开。它到像又是双射，各局部光滑逆在交叠处由唯一性相同，便拼成全局光滑逆。
\end{proof}

\cref{g1:fig:free-action-product}显示这里的两个参数：群参数沿轨道变化，\(S\) 上的参数在不同轨道之间变化。图中画出的只是 \(G\) 的一小块；定理中的乘积却包含整个 \(G\)。因此它并不声称轨道本身像直线，也不要求 \(G\) 连通或紧。

\begin{figure}[htbp]
	\centering
	% !TeX root = ../Book1.tex
\begin{tikzpicture}[>=stealth,font=\small]
	\fill[blue!4] (-2,-1.2) rectangle (2,1.2);
	\foreach \y/\colour in {-0.8/blue!70!black,0/red!70!black,0.8/green!50!black}{
		\draw[thick,\colour,->] (-1.8,\y)--(1.8,\y);
		\fill[\colour] (0,\y) circle (2pt);
		\fill[\colour] (4,\y) circle (2pt);
	}
	\draw[thick] (0,-1.1)--(0,1.1) node[above] {\(\{e\}\times S\)};
	\draw[thick] (4,-1.1)--(4,1.1) node[above] {\(S\cong\pi(GS)\)};
	\draw[->] (2.3,0.2)--(3.4,0.2) node[midway,above] {\(\pi\)};
	\node at (0,-1.6) {\(G\times S\cong GS\)};
	\node at (4,-1.6) {横向坐标};
\end{tikzpicture}

	\caption[自由固有作用的局部乘积]{局部乘积 \(G\times S\cong GS\) 中的轨道方向与横向方向。每条水平线表示一条轨道的一部分，竖线表示横截面 \(\{e\}\times S\)；商映射忘掉群参数，只保留横向坐标。图中直线仅表示乘积坐标，不表示原流形中的轨道形状。}
	\label{g1:fig:free-action-product}
\end{figure}

\subsection{商的 Hausdorff 性和光滑图册}
\label{g1:subsec:1702:03}

给轨道集 \(M/G\) 赋予商拓扑，记 \(\pi(p)=[p]=Gp\)。对任意开集 \(O\subseteq M\)，有
\[
	\pi^{-1}(\pi(O))=GO=\bigcup_{g\in G}gO,
\]
它是开集，故 \(\pi\) 为开映射。这一点对任意连续群作用都成立；要得到流形，还须证明不同轨道能够分开，并为这些轨道配上相容的光滑坐标。

\begin{theorem}[自由固有作用的商流形]\label{g1:thm:free-proper-quotient}
	设 \(G\) 为 \(r\) 维实 Lie 群，\(M\) 为 \(m\) 维无边界光滑流形，两者均为 Hausdorff、第二可数。若 \(G\) 自由且固有地光滑作用于 \(M\)，则商拓扑下的 \(M/G\) 具有唯一的 \(m-r\) 维无边界光滑流形结构，使投影
	\[
		\pi:M\longrightarrow M/G
	\]
	成为光滑满射淹没。在\cref{g1:lem:free-proper-local-product}的每个局部乘积 \(G\times S\cong GS\) 中，\(\pi\) 对应第二因子投影，且 \(\pi|_S:S\to\pi(GS)\) 是微分同胚。
\end{theorem}
\begin{proof}
	先核验商空间的分离性与可数性。若 \([p]\ne[q]\)，则 \((q,p)\) 不在闭关系集 \(R\) 中。由\cref{g1:prop:proper-action-orbits}，可取开邻域 \(U\ni p\)、\(V\ni q\)，使 \((V\times U)\cap R=\varnothing\)。这就是说没有轨道同时经过 \(U\) 和 \(V\)，所以 \(\pi(U)\)、\(\pi(V)\) 是两个不交的开邻域，分别包含 \([p]\)、\([q]\)。故商空间 Hausdorff。取 \(M\) 的可数开基 \((B_j)\)，开集 \(\pi(B_j)\) 构成商空间的可数基：对 \([p]\in W\) 且 \(W\) 开，选 \(p\) 的基邻域 \(B_j\subseteq\pi^{-1}(W)\)，即可得到 \([p]\in\pi(B_j)\subseteq W\)。

	对局部横截面 \(S\)，记 \(W=\pi(GS)\)。\(GS\) 开且饱和，故 \(W\) 开；每条轨道与 \(S\) 交于唯一点，使 \(\pi|_S:S\to W\) 成为连续双射。它还开：若 \(A\subseteq S\) 开，则 \(G\times A\) 在 \(G\times S\) 中开，局部乘积微分同胚把它送到 \(GA\)，这是 \(M\) 的开集，所以 \(\pi(A)=\pi(GA)\) 在商中开。因此 \(\pi|_S\) 是同胚。取 \(S\) 上的光滑坐标 \(u:S\to\mathbb R^{m-r}\) 到一个开集，便得到 \(W\) 上的坐标 \(u\circ(\pi|_S)^{-1}\)。

	检验交叠时，设另一个横截面为 \(S'\)，局部乘积逆映射为 \(\Phi'^{-1}:GS'\to G\times S'\)。若 \(s\in S\cap GS'\)，唯一写成 \(s=g s'\)，其中 \(s'\in S'\)。用两张横截面表示同一轨道时，相应的转换正是
	\[
		S\cap GS'\longrightarrow S',\qquad
		s\longmapsto\operatorname{pr}_{S'}\bigl(\Phi'^{-1}(s)\bigr).
	\]
	定义域在 \(S\) 中开，右侧是光滑映射的复合，因此转换光滑；交换 \(S,S'\) 后得到光滑逆。于是上述坐标构成相容光滑图册，维数为 \(m-r\)，且没有边界。

	在乘积坐标中，\(\pi(gs)=\pi(s)\)，所以 \(\pi\) 就是 \((g,s)\mapsto s\)。它光滑且微分满射，因而为淹没；定义又使它满射。这些坐标还给出每个商点附近的光滑截面 \(\sigma:W\to M\)，即 \(\sigma([s])=s\)。

	最后证明唯一性。设同一商拓扑上还有另一份光滑结构，使同一个投影 \(\pi'\) 为光滑满射淹没。从刚构造的商到另一份商的恒等映射，在 \(W\) 上可写为 \(\pi'\circ\sigma\)，故光滑。反方向中，\cref{g1:thm:immersion-submersion-normal-form}说明 \(\pi'\) 的局部表达为投影；把被忘掉的坐标取常值，就得到目标点附近的光滑截面 \(\tau\)。从另一份商到本商的恒等映射在这片邻域上为 \(\pi\circ\tau\)，也光滑。两个恒等映射互逆，所以两份光滑结构相同。
\end{proof}

自由性使轨道方向始终具有维数 \(\dim G\)，固有性则既排除横截面的远处返回，也保证商的 Hausdorff 性。证明中第二可数性同样保留下来，不能只凭“每点附近看起来像 Euclid 空间”就称所得空间为本书意义的流形。

\begin{example}[正比例缩放的商]\label{g1:exm:positive-scaling-quotient}
	乘法 Lie 群 \((0,\infty)\) 作用于 \(\mathbb R^n\setminus\{0\}\)、\(n\ge1\)，作用为 \(a\cdot x=ax\)。它自由。若 \(x_j\to x\ne0\)、\(a_jx_j\to y\ne0\)，则
	\[
		a_j=\frac{\lVert a_jx_j\rVert}{\lVert x_j\rVert}
		\longrightarrow\frac{\lVert y\rVert}{\lVert x\rVert}>0,
	\]
	故作用固有。每条轨道恰与单位球面 \(S^{n-1}\) 相交一次，并且
	\[
		(0,\infty)\times S^{n-1}\longrightarrow\mathbb R^n\setminus\{0\},
		\qquad(a,u)\longmapsto au
	\]
	的光滑逆为 \(x\mapsto(\lVert x\rVert,x/\lVert x\rVert)\)。所以这里甚至有全局乘积，商就是 \(S^{n-1}\)，商投影对应 \(x\mapsto x/\lVert x\rVert\)。在 \(n=1\) 时，球面 \(S^0\) 为两个点；正负半轴分别是一条轨道，这也与零维商一致。
\end{example}

\subsection{商映射的泛性质}
\label{g1:subsec:1702:04}

拓扑商的泛性质见\cref{g1:prop:quotient-descent}：在每条轨道上取相同值的连续映射，才能成为商上的连续映射。现在局部横截面使同一结论提升到光滑范畴。实际使用时，这通常比直接计算商坐标更方便。

\begin{proposition}[光滑映射的下降]\label{g1:prop:quotient-smooth-universal-property}
	设有限维实 Lie 群 \(G\) 自由且固有地光滑作用于有限维无边界光滑流形 \(M\)，并设 \(N\) 为有限维无边界光滑流形；所有流形均为 Hausdorff、第二可数。给 \(M/G\) 以上一定理的光滑结构。则集合映射 \(h:M/G\to N\) 光滑，当且仅当 \(h\circ\pi:M\to N\) 光滑。

	特别地，光滑映射 \(f:M\to N\) 满足 \(f(gp)=f(p)\) 对所有 \(g,p\) 成立，当且仅当存在唯一光滑映射 \(\bar f:M/G\to N\)，使 \(f=\bar f\circ\pi\)。
\end{proposition}
\begin{proof}
	若 \(h\) 光滑，复合 \(h\circ\pi\) 光滑。反方向中，取商空间任意点附近的光滑截面 \(\sigma:W\to M\)，则
	\[
		h|_W=(h\circ\pi)\circ\sigma,
	\]
	所以它光滑。这样的 \(W\) 覆盖整个商，故 \(h\) 光滑；证明甚至不必预先假定 \(h\) 连续。

	若 \(f\) 在轨道上常值，定义 \(\bar f([p])=f(p)\)。轨道常值性保证定义不依赖代表元，满射性保证存在且唯一，刚证的光滑性判据保证 \(\bar f\) 光滑。反过来，任何复合 \(\bar f\circ\pi\) 在轨道上都取常值。
\end{proof}

这一泛性质可以画成
\[
	\begin{tikzcd}[column sep=large,row sep=large]
		M \arrow[r,"f\text{ 光滑且不变}"] \arrow[d,two heads,"\pi\text{ 淹没}"'] & N \\
		M/G \arrow[ur,dashed,"\bar f\text{ 唯一光滑}"'] &
	\end{tikzcd}
\]
图中的 \(N\) 是任意测试流形，不承担给定的 \(G\) 作用。若目标也带作用，等变映射未必沿源轨道取常值；要下降到轨道商，通常应同时把目标送到它的轨道商，不能直接套用这里的不变性条件。

\subsection{缺少自由或固有条件的反例}
\label{g1:subsec:1702:05}

\begin{example}[自由作用的非 Hausdorff 商]\label{g1:exm:irrational-flow-quotient}
	固定无理数 \(\alpha\)，令加法群 \(\mathbb R\) 作用于 \(T^2=\mathbb R^2/\mathbb Z^2\)：
	\[
		t\cdot[x,y]=[x+t,y+\alpha t].
	\]
	若某个 \(t\) 固定一点，则 \(t\in\mathbb Z\)、\(\alpha t\in\mathbb Z\)，所以 \(t=0\)，作用自由。\cref{g1:exm:irrational-torus-subgroup}已证明单位轨道是环面的真稠密子群；其他轨道是它的平移，也都稠密。取非空开集 \(O\subseteq T^2\)，每条轨道都与 \(O\) 相交，所以其饱和集 \(\mathbb R O\) 为整个环面。于是轨道商只有空集和全空间两个开集，而又有不止一条轨道，因此不是 Hausdorff 空间。

	这里固有性明确失败：紧集 \(K=T^2\) 的搬运集 \(T(K,K)=\mathbb R\) 非紧。每条轨道本身仍有一维浸入流形结构；失败的是把所有轨道合并以后，无法给它们彼此分离的邻域。
\end{example}

\begin{example}[旋转商的端点]\label{g1:exm:rotation-half-line-quotient}
	圆周群 \(S^1\) 通过复数乘法作用于 \(\mathbb C\cong\mathbb R^2\)。群紧，所以作用固有；原点的稳定子是整个 \(S^1\)，其他点的稳定子平凡，因此作用不自由。轨道由半径确定，\([z]\mapsto|z|\) 给出
	\[
		\mathbb C/S^1\cong[0,\infty).
	\]
	这是同胚：模长连续且在轨道上常值，故诱导连续映射；其逆 \(r\mapsto[r]\) 由正实半轴的连续包含给出。端点附近删去端点以后仍连通，而开区间中删去内点会出现两个分支，所以该商不是一维无边界流形。

	即使允许边界，商投影也不会满足前一定理的淹没结论。常用的光滑不变量 \(u(z)=|z|^2\) 在原点的微分为零。更一般地，若实值光滑函数 \(f\) 对旋转不变，则 \(f(-z)=f(z)\)，在原点求导得到 \(\mathrm{d}f_0=-\mathrm{d}f_0\)，因此 \(\mathrm{d}f_0=0\)。任何通过一维商坐标描述的光滑投影，在原点都不可能有秩 \(1\)。
\end{example}

非自由性也不意味着商一定出现坏点。例如 \(S^1\) 在任意流形 \(M\) 上的平凡作用是固有而不自由的，轨道商仍是 \(M\)，投影就是恒等映射。改变的是“轨道维数等于 \(\dim G\)”这一前提，故不能使用 \(\dim M-\dim G\) 的维数公式。下一步研究非自由作用时，需要保留稳定子及其横向作用，才能区分这种无害的冗余对称与旋转例中的端点。

% 纲要条目排为编号小节标题，并进入目录与书签；说明段仍保留于纲要注释中。
